% =============================================================================
%  STUDY COMPANION — Course Overview and NLP Applications
%  ZG519 NLP Applications, Session 1
%  Built with make-lecture-kit v3.1.0
% =============================================================================

\documentclass[a4paper,11pt]{article}

\usepackage[T1]{fontenc}
\usepackage[utf8]{inputenc}
\IfFileExists{mathpazo.sty}{\usepackage[sc]{mathpazo}\linespread{1.06}}{\usepackage{lmodern}}
\IfFileExists{helvet.sty}{\usepackage[scaled=0.9]{helvet}}{}
\renewcommand{\ttdefault}{lmtt}
\usepackage{microtype}

\usepackage[a4paper,top=26mm,bottom=26mm,left=27mm,right=27mm,
            headheight=14pt,headsep=12mm,footskip=12mm]{geometry}

\usepackage{amsmath,amssymb}

\usepackage{graphicx}
\usepackage{booktabs}
\usepackage{array}
\usepackage[table,svgnames]{xcolor}
\usepackage{caption}
\usepackage{float}
\usepackage[section]{placeins}
\usepackage{needspace}

\usepackage[most]{tcolorbox}
\tcbuselibrary{skins,breakable}

\usepackage{tikz}
\usepackage{pifont}

\usepackage{enumitem}
\setlist{leftmargin=1.5em,topsep=4pt,itemsep=3pt}

\newlist{steps}{enumerate}{1}
\setlist[steps]{label={\sffamily\small\bfseries\color{accent}Step \arabic*},
  leftmargin=4.2em, labelwidth=3.6em, labelsep=0.6em, labelindent=0pt,
  align=left, topsep=6pt, itemsep=5pt}
\usepackage{parskip}
\setlength{\parskip}{0.55em plus 0.1em}

\usepackage{fancyhdr}
\usepackage{titlesec}

\usepackage[hidelinks]{hyperref}
\usepackage{xurl}
\hypersetup{breaklinks=true}

\definecolor{accent}{HTML}{1F4E79}
\definecolor{inktext}{HTML}{1F2328}
\definecolor{mutedtext}{HTML}{6B7280}
\definecolor{hairline}{HTML}{D0D5DD}
\definecolor{wash}{HTML}{F5F6F8}
\colorlet{navy}{accent}      \colorlet{navyrule}{hairline}
\colorlet{intuFrame}{accent} \colorlet{everyFrame}{accent} \colorlet{workFrame}{accent}
\colorlet{watchFrame}{inktext} \colorlet{keyFrame}{accent} \colorlet{waitFrame}{accent}

\color{inktext}
\hypersetup{linkcolor=accent,urlcolor=accent,citecolor=accent}

\newcommand{\SeriesName}{ZG519 NLP Applications}
\newcommand{\TopicTitle}{Course Overview \& NLP Applications}

\pagestyle{fancy}
\fancyhf{}
\fancyhead[L]{\sffamily\scriptsize\color{mutedtext}\MakeUppercase{\SeriesName}}
\fancyhead[R]{\sffamily\scriptsize\color{mutedtext}\TopicTitle}
\fancyfoot[C]{\sffamily\footnotesize\color{mutedtext}\thepage}
\renewcommand{\headrulewidth}{0pt}
\renewcommand{\footrulewidth}{0pt}
\fancypagestyle{plain}{\fancyhf{}\fancyfoot[C]{\sffamily\footnotesize\color{mutedtext}\thepage}}

\titleformat{\section}
  {\sffamily\Large\bfseries\color{inktext}}
  {\color{accent}\thesection}{0.7em}{}
\titleformat{\subsection}
  {\sffamily\large\bfseries\color{inktext}}
  {\color{accent}\thesubsection}{0.6em}{}
\titlespacing*{\section}{0pt}{26pt}{6pt}
\titlespacing*{\subsection}{0pt}{14pt}{4pt}
\newcommand{\sectionbreak}{\needspace{14\baselineskip}}

\newcommand{\secsub}[1]{%
  \noindent{\large\itshape\color{mutedtext}#1}\par
  \vspace{2pt}{\color{hairline}\rule{\linewidth}{0.4pt}}\par\medskip}

\newcommand{\flipanswer}[1]{\par\noindent\hfill\rotatebox[origin=c]{180}{\parbox{0.9\linewidth}{\small\color{mutedtext}#1}}\par\medskip}

\tcbset{
  quiet/.style 2 args={
    enhanced, breakable, sharp corners,
    colback=#2, colframe=#2, boxrule=0pt,
    borderline west={1.6pt}{0pt}{#1},
    left=12pt, right=10pt, top=7pt, bottom=7pt,
    before skip=12pt, after skip=12pt,
    fontupper=\normalsize, coltitle=#1,
    fonttitle=\sffamily\scriptsize\bfseries,
    attach title to upper={\par\smallskip},
  }
}
\newcommand{\calloutlabel}[1]{\MakeUppercase{#1}}

\newtcolorbox{intuition}{quiet={accent}{white},
  title={\calloutlabel{The intuition}}}
\newtcolorbox{everyday}{quiet={accent}{white},
  title={\calloutlabel{Everyday picture}}}
\newtcolorbox{worked}[1]{quiet={accent}{wash},
  title={\calloutlabel{Worked example}\ \textperiodcentered\ {\normalfont\itshape\small\color{inktext}#1}}}
\newtcolorbox{watchout}{quiet={inktext}{white},
  title={\calloutlabel{Watch out}}}
\newtcolorbox{waitwhat}{quiet={accent}{white},
  title={\calloutlabel{Wait --- really?}}}
\newtcolorbox{keytake}{quiet={accent}{white},
  borderline north={0.4pt}{0pt}{hairline},
  borderline south={0.4pt}{0pt}{hairline},
  title={\upshape\calloutlabel{Key takeaway}}, fontupper=\normalsize\itshape}

\newcommand{\titlebanner}[4]{%
  \begingroup\raggedright
  \vspace*{4mm}
  {\sffamily\footnotesize\bfseries\color{accent}\MakeUppercase{#1}}\par
  \vspace{7pt}
  {\sffamily\bfseries\fontsize{30}{35}\selectfont\color{inktext}#2\par}
  \vspace{7pt}
  {\large\color{mutedtext}#3\par}
  \vspace{12pt}
  {\color{hairline}\rule{\linewidth}{0.5pt}}\par
  \vspace{3pt}
  {\sffamily\scriptsize\color{mutedtext}\MakeUppercase{#4}}\par
  \vspace{10pt}
  \endgroup
}

\captionsetup[figure]{font={small,color=mutedtext},labelfont={bf,sf,color=inktext},
  labelsep=period,justification=raggedright,singlelinecheck=false,skip=6pt}
\setlength{\textfloatsep}{16pt plus 4pt minus 4pt}
\setlength{\intextsep}{14pt plus 4pt minus 4pt}
\renewcommand{\topfraction}{0.8}\renewcommand{\bottomfraction}{0.6}
\renewcommand{\textfraction}{0.15}\renewcommand{\floatpagefraction}{0.75}
\newcommand{\housefig}[3][0.86]{%
  \begin{figure}[!htbp]
    \centering
    \includegraphics[width=#1\linewidth]{#2}%
    \caption{#3}
  \end{figure}%
}

\renewcommand{\arraystretch}{1.18}
\newcommand{\vocab}[1]{\textbf{#1}}

\begin{document}

\thispagestyle{fancy}

\titlebanner
  {\SeriesName}
  {Course Overview \& NLP Applications}
  {Ten application areas, easiest to hardest, and why the easy ones still run in production}
  {Built from: the Session 1 slides, course handout and recording · Every example worked out in full}

\noindent\textbf{How to use this companion.}
This companion explains Session 1 from scratch, in plain English. Every idea
starts with an everyday picture. Then the mechanism is built up one step at a
time. Every example is worked out in full. Small labels mark the boxes.
\textsf{\small\textbf{THE INTUITION}} states the core idea in one breath.
\textsf{\small\textbf{EVERYDAY PICTURE}} ties it to real life before any
mechanism. \textsf{\small\textbf{WORKED EXAMPLE}}, on a grey background,
walks a real task through with no skipped steps. \textsf{\small\textbf{WATCH
OUT}} flags the common mistake. \textsf{\small\textbf{KEY TAKEAWAY}} is the
one sentence to remember. Read it in order; each section leans on the one
before.

\medskip

% =============================================================================
\section{Try it first: what does one messy sentence need?}
\secsub{Before any course map, feel how many separate NLP steps a single
real message already demands.}

Here is a WhatsApp message. It is exactly the kind an NLP system meets daily:

\begin{quote}
\itshape
``Bro, idhu kashtam (this is hard) but deadline tmrw, pls extend!!''
\end{quote}

Before reading on, picture handing this sentence straight to a translation
tool, with no other processing. Ask three questions. Does the tool know
\emph{which words are English and which are Tamil}? Does it know the writer
is \emph{stressed, asking for a favour}, not stating a neutral fact? Does it
even know ``tmrw'' means ``tomorrow''? A raw translation tool assumes clean,
single-language, correctly-spelled input. This sentence is none of those
three things, and real text almost never is.

That gap names three separate jobs, each its own application area this
semester. First, \vocab{spelling and grammar correction} turns ``tmrw'' into
``tomorrow'' before anything downstream can use it cleanly. Second,
\vocab{language identification} has to spot that ``idhu kashtam'' is Tamil,
mixed into an English sentence --- called \vocab{code-switching}. Third,
\vocab{sentiment analysis} has to read the stressed, pleading tone from
``pls extend!!'', not just the dictionary meaning of the words. One short
message, three applications, before any translation or question-answering
even starts. That ordering --- clean first, understand next, act last --- is
the backbone of this entire course.

% =============================================================================
\section{Sixteen sessions, one map}
\secsub{Every session this semester, and exactly how the two exams are
weighted, straight from the official handout.}

\begin{table}[htbp]
\centering\small
\begin{tabular}{@{}p{0.9cm}p{4.0cm}p{6.3cm}@{}}
\toprule
\textbf{\#} & \textbf{Topic} & \textbf{Sub-topics} \\
\midrule
1  & Introduction & Overview of interesting NLP applications \\
2  & Spellcheck \& grammar check & Techniques for automatic checkers \\
3  & Retrieval-Augmented Generation & Na\"ive RAG, Graph RAG, Agentic RAG \\
4  & Question answering & Agentic AI + Gen AI techniques in QA \\
5  & Conversational AI & Hybrid systems: LLM + RAG + Agentic AI \\
6  & Sentiment analysis & Methods; neural networks for sentiment \\
7  & Privacy \& ethics in NLP & Data, model, and deployment issues \\
8  & Pre-midsem review & Review of Sessions 1--7 \\
9  & Machine translation & Statistical approaches \\
10 & Machine translation & Neural approaches \\
11 & Machine translation & Indic language translation \\
12 & Information extraction & Named entity recognition \\
13 & Information extraction & Relation extraction \\
14 & Information extraction & Events and time extraction \\
15 & Scalable \& efficient NLP & Compression, inference, deployment \\
16 & Post-midsem review & Review of Sessions 9--15 \\
\bottomrule
\end{tabular}
\caption*{\footnotesize All sixteen contact sessions, compressed.}
\end{table}

\begin{table}[htbp]
\centering\small
\begin{tabular}{@{}p{4.6cm}p{1.4cm}p{4.2cm}@{}}
\toprule
\textbf{Component} & \textbf{Weight} & \textbf{Notes} \\
\midrule
2 quizzes (pre- \& post-midsem) & 10\% & Open on mobile; mark dates the moment they post \\
Programming assignment(s) & 25\% & Groups of 4, end-to-end applications \\
Situated learning & 5\% & Apply one application to your own workplace \\
Mid-semester exam & 30\% & Closed book, Sessions 1--8 only \\
Comprehensive final exam & 35\% & Open book, all topics \\
\bottomrule
\end{tabular}
\caption*{\footnotesize How the course is graded.}
\end{table}

\begin{watchout}
The exam that rewards memorisation is the mid-sem, not the final. The
mid-sem is closed book and covers only Sessions 1--8: spelling, RAG basics,
QA, conversational AI, sentiment, and ethics. Learn those as you go. The
comprehensive final is open book, so it rewards fluency with the tools, not
recall.
\end{watchout}

% =============================================================================
\section{Ten applications, one difficulty ladder}
\secsub{The professor orders every application by difficulty on purpose.
Knowing when a cheap classical method beats an expensive agent is itself
an exam skill.}

A fair question up front. Large language models can do almost anything
with text, so why does this course still teach ten separate,
old-fashioned-sounding applications? Because production systems are
\vocab{hybrid}. The simple method is often cheaper and faster. It is also
often just as accurate as the fancy one. Knowing which to reach for is
the actual skill this course builds.

\begin{intuition}
Every later era of a technology keeps using the earlier one wherever it is
good enough. A spam filter does not need an agent. A multi-hop research
question might. The ladder below is a map of \emph{when} to reach for more
machinery, not a claim that newer always wins.
\end{intuition}

\housefig{figures/fig_difficulty.png}{Ten applications, ranked by how much
machinery (and cost) they typically need, from a near-perfect classical
method to a full agent loop.}

\paragraph{1 $\cdot$ Spell \& grammar check --- the baseline.}
Classic statistical methods (Naive-Bayes-style) already get this close to
100\% right, no large language model required. It still matters: it is the
first filter in front of every later application. A clean audio recording
has to exist, the same way, before anyone can transcribe it.

\paragraph{2 $\cdot$ Language identification \& code-switching.}
You must know the language before you can translate, tag, or search it. The
genuinely hard case is \vocab{code-switching}: one sentence mixing two
languages, exactly like the WhatsApp message in Section 1.

\begin{worked}{tagging the code-switched message word by word}
\begin{steps}
  \item Read each word and ask: which language's vocabulary does this word
    belong to?
  \item ``Bro,'', ``but'', ``deadline'', ``tmrw'' --- English spellings and
    grammar, even where ``tmrw'' is a shortened spelling.
  \item ``idhu'', ``kashtam'' --- Tamil words (``this'', ``hard''), written
    in English letters, a common style called transliteration.
  \item Tag each token EN or TA. The output is not one label for the whole
    sentence. It is a label \emph{per word}. That is exactly why
    code-switching breaks single-language detectors: they expect one
    answer for the whole input.
\end{steps}
\end{worked}

\housefig{figures/fig_codeswitch.png}{Language identification on a
code-switched sentence: each word gets its own EN/TA tag, not one tag for
the whole sentence.}

\paragraph{3 $\cdot$ Sentiment analysis.}
Classifying text as positive, negative, or neutral --- a movie review, a
product review, a political statement. The signal comes from word choice,
punctuation (``!!''), and emphasis, not just the dictionary meaning of each
word.

\paragraph{4 $\cdot$ Summarization: extractive versus abstractive.}
\vocab{Extractive} summarization copies real sentences straight from the
source. \vocab{Abstractive} summarization writes new sentences that capture
the gist, in the system's own words. It is much harder, because it
requires actually following the content, not just picking good excerpts.

\paragraph{5 $\cdot$ Information extraction --- entities, relations, events.}
Covered in full depth in Section 4 below: turning a plain sentence into
structured facts a database can store.

\paragraph{6 $\cdot$ Machine translation.}
A headline topic this run, with three dedicated sessions. Statistical MT
comes first, then neural MT, then Indic-language MT specifically, tied to
national efforts like AI4Bharat and Sarvam AI.

\paragraph{7 $\cdot$ Question answering \& conversational AI.}
A huge open-source and commercial tooling ecosystem lives here. The course
builds these as hybrid systems --- LLM plus RAG plus Agentic AI. A bare
LLM alone hallucinates on specific factual questions.

\paragraph{8 $\cdot$ RAG --- na\"ive, graph, and agentic.}
The hardest application area in the course. Covered in full depth in
Section 5 below.

\paragraph{9 $\cdot$ Agentic AI \& guardrails.}
Agentic systems carry a real cost. Production teams do not reach for them
by default; they escalate only where simpler methods fail. Whenever agents
are used, guardrails and a human in the loop are treated as non-negotiable,
not optional extras.

\paragraph{10 $\cdot$ Scalable \& efficient NLP systems.}
New to this run. It was added because a method that works beautifully on a
small test set can fail in production. Latency, scale, and cost problems
often never show up in a notebook. Session 15 covers model compression and
inference optimisation directly.

\begin{keytake}
Ten applications, one ladder: near-perfect classical methods at the easy
end, a full agent loop with guardrails at the hard end. The skill this
course teaches is not ``always use the newest method.'' It is knowing
which rung of the ladder a given problem actually needs.
\end{keytake}

% =============================================================================
\section{Worked example: turning a sentence into structured facts}
\secsub{Information extraction has three layers: who is mentioned, what did
they do to whom, and when.}

\vocab{Named entity recognition} (NER) finds proper nouns --- people,
places, organisations --- in running text. \vocab{Relation extraction}
then connects two entities it found. \vocab{Event extraction} adds the
``when''. All three lean on knowing the \vocab{part of speech} of each
word, since named entities typically begin with a proper noun.

\begin{everyday}
This is exactly what IMDb does automatically to a movie review. It tags
the cast, the director, and the studio it mentions --- entities. It works
out who directed whom --- a relation. It notes the release date the review
refers to --- an event. All of this comes from one paragraph of free text.
\end{everyday}

\begin{worked}{extracting structured facts from one sentence}
Sentence: ``Google acquired DeepMind in 2014 for Rs.~4{,}000 crore.''
\begin{steps}
  \item \textbf{Tokenise \& tag part of speech.} Split into words; mark
    ``Google'' and ``DeepMind'' as proper nouns, ``acquired'' as a verb,
    ``2014'' and ``4{,}000'' as numbers.
  \item \textbf{Named entity recognition.} Label each proper noun and
    number span with its entity type: \texttt{Google = ORG},
    \texttt{DeepMind = ORG}, \texttt{2014 = DATE},
    \texttt{Rs.~4{,}000 crore = MONEY}.
  \item \textbf{Relation extraction.} The verb ``acquired'' links the two
    ORG entities in a fixed direction: \texttt{acquired(Google, DeepMind)}.
    Google is the buyer. DeepMind is the one bought, not the reverse.
  \item \textbf{Event extraction.} Attach the DATE and MONEY entities to
    that relation as the event's details. The result: an acquisition
    event, subject Google, object DeepMind, time 2014, cost Rs.~4{,}000
    crore.
\end{steps}
The output is no longer a sentence. It is a row a database can store:
one event, four typed fields, ready to query.
\end{worked}

\housefig{figures/fig_ner.png}{Named entity tags on the acquisition
sentence: ORG, DATE, and MONEY spans, each found independently before being
linked into one event.}

\begin{watchout}
NER alone does not tell you \emph{who acquired whom}. ``Google acquired
DeepMind'' and ``DeepMind acquired Google'' contain the exact same two
entities. Only relation extraction, which reads the verb and its direction,
tells the two apart. Entities, relations, and events are three separate
jobs, not one.
\end{watchout}

\begin{keytake}
Information extraction builds structured facts in layers. Find the
entities first. Link pairs of them into relations. Then attach time and
detail to make a full event. Each layer depends on the one before it
being right.
\end{keytake}

% =============================================================================
\section{Worked example: the RAG ladder --- na\"ive, graph, agentic}
\secsub{Three ways to answer one question, each fixing the previous one's
blind spot on multi-step questions.}

\vocab{Retrieval-Augmented Generation} (RAG) answers a question in two
steps. First it retrieves relevant text. Then it generates an answer
grounded in that text, instead of relying only on what the model
memorised during training.

\begin{everyday}
Na\"ive RAG is a friend who looks up one relevant fact before answering your
question. Graph RAG is a friend who also knows how facts connect to each
other. They can answer ``how is X related to Y'' by walking a web of
relationships, not just reading isolated snippets. Agentic RAG is that same
friend deciding, mid-conversation, ``actually, I should look that up again
before I answer this follow-up.'' Retrieval becomes a decision for them to
make, not a fixed, one-time step.
\end{everyday}

\begin{worked}{answering a two-hop question three different ways}
Question: ``What company did DeepMind's co-founder previously work at
before DeepMind was acquired by Google?''
\begin{steps}
  \item \textbf{Na\"ive RAG.} Retrieve the single text chunk most similar to
    the question --- likely a chunk about the Google–DeepMind acquisition.
    That chunk may never mention the co-founder's earlier job. The system
    answers confidently from an incomplete chunk, or admits it does not
    know.
  \item \textbf{Graph RAG.} The underlying knowledge graph stores typed
    edges: \texttt{founded(Person, DeepMind)}, \texttt{acquired(Google,
    DeepMind)}, \texttt{previously\_worked\_at(Person, Company)}. Following
    the \texttt{Person} node across all three edges answers the question
    in one lookup. The connection between facts is explicit in the graph,
    not buried in prose.
  \item \textbf{Agentic RAG.} With plain text chunks and no graph, the
    system reasons for itself. ``The first chunk names the co-founder, but
    not their earlier job. I need a second, different retrieval.'' It
    retrieves again, this time searching for that person's name plus
    ``previously worked at''. Two retrieval rounds, chosen by the system
    itself, replace the one graph lookup.
\end{steps}
All three reach the same correct answer here. Na\"ive RAG is the one that
quietly fails on real two-hop questions like this one. It never gets a
second chance to look further.
\end{worked}

\housefig{figures/fig_agentic_rag.png}{Agentic RAG as a loop: the system
decides, after each retrieval, whether it has enough evidence yet or needs
to look again.}

\begin{waitwhat}
Agentic RAG does not use a smarter retrieval model than na\"ive RAG. The
retrieval step is identical. What changed is that retrieval now sits
\emph{inside a loop} the system controls. That is the same
reason-act-observe pattern that turns a chatbot into an agent elsewhere
in this course.
\end{waitwhat}

\begin{keytake}
Na\"ive RAG retrieves once, fixed. Graph RAG retrieves across explicit
connections between facts. Agentic RAG retrieves as many times as the
system itself decides it needs to. Multi-hop questions are where na\"ive
RAG's one-shot limit shows up first.
\end{keytake}

\noindent\textbf{Where it's used.}
This same ladder explains a common failure. A company's internal search
bot often answers simple FAQs correctly. It can still fail on ``compare
last year's and this year's policy on X.'' That is a two-hop question, and
plain, single-pass RAG was never built to answer it.

% =============================================================================
\section*{Self-test}
\addcontentsline{toc}{section}{Self-test}

\noindent\textbf{1.} Why does spelling and grammar correction run before
almost every other NLP application, not after?

\flipanswer{Every later step --- language ID, sentiment, NER, translation
--- assumes reasonably clean input text. An uncorrected ``tmrw'' or a
mis-split sentence makes every downstream step less accurate.}

\noindent\textbf{2.} Why do single-language detectors fail on
code-switched text?

\flipanswer{They are built to output one label for the whole input. A
code-switched sentence genuinely needs a different label per word, so a
single whole-sentence answer is wrong by construction.}

\noindent\textbf{3.} In ``Google acquired DeepMind'', what does relation
extraction add that named entity recognition alone cannot tell you?

\flipanswer{Direction. NER finds that Google and DeepMind are both
organisations, but only the verb-driven relation extraction step tells you
\emph{which one acquired which}.}

\noindent\textbf{4.} Why does na\"ive RAG struggle on a two-hop question
like the DeepMind co-founder example?

\flipanswer{It retrieves only once. If the single best-matching chunk
does not contain every fact the answer needs, there is no second chance
to look further. Graph RAG has explicit connections instead, and agentic
RAG can simply retrieve again.}

\noindent\textbf{5.} Why does the mid-semester exam reward memorisation
more than the comprehensive final does?

\flipanswer{The mid-sem is closed book and covers only Sessions 1--8. The
comprehensive final is open book and covers everything, so it rewards
fluency using the material, not recalling it from memory.}

% =============================================================================
\section*{Pocket card: Session 1 in one page}
\addcontentsline{toc}{section}{Pocket card}

\begin{keytake}
\textbf{The whole session, compressed.}
\begin{itemize}[leftmargin=1.3em,topsep=4pt,itemsep=4pt]
  \item One messy real sentence already needs three separate applications
    --- \textbf{spelling correction}, \textbf{language ID}
    (code-switching), and \textbf{sentiment}. Clean first, understand
    next, act last.
  \item Sixteen sessions, two exams: mid-sem is \textbf{closed book},
    Sessions 1--8 only; the final is \textbf{open book}, all topics.
  \item Ten applications form one \textbf{difficulty ladder}, from
    near-perfect classical spell-check to a full agent loop with
    guardrails. Production systems are \textbf{hybrid}: use the cheapest
    method that works.
  \item \textbf{Information extraction} builds structured facts in layers:
    entities $\to$ relations $\to$ events. Each layer needs the one before
    it to be correct.
  \item The \textbf{RAG ladder}: na\"ive RAG retrieves once. Graph RAG
    follows explicit connections between facts. Agentic RAG decides, on
    its own, whether to retrieve again. Multi-hop questions expose
    na\"ive RAG's one-shot limit first.
\end{itemize}
\end{keytake}

% =============================================================================
\section*{Glossary \& symbols}
\addcontentsline{toc}{section}{Glossary \& symbols}

\noindent\textbf{Key terms.}
\begin{description}[leftmargin=9em,style=nextline,font=\normalfont\bfseries\color{navy}]
  \item[code-switching] one sentence or message that mixes two or more
    languages, often within a single line.
  \item[extractive summary] a summary built only from sentences copied
    verbatim from the source text.
  \item[abstractive summary] a summary written in new words that capture
    the gist, without copying sentences.
  \item[named entity recognition (NER)] finding spans of text that name a
    person, place, organisation, date, or amount.
  \item[relation extraction] linking two entities already found by NER
    into a typed, directional connection (who did what to whom).
  \item[event extraction] attaching time and other detail to a relation to
    describe a full real-world event.
  \item[RAG] Retrieval-Augmented Generation: retrieve relevant text, then
    generate an answer grounded in it.
  \item[graph RAG] RAG where retrieval walks explicit, typed connections
    between facts in a knowledge graph, not flat text chunks.
  \item[agentic RAG] RAG where the system itself decides when and what to
    retrieve, possibly retrieving more than once.
  \item[hybrid system] a production system that deliberately mixes cheap
    classical methods with expensive LLM/agentic methods, by task.
  \item[guardrail] a hard limit placed on an agent's actions, to keep
    autonomy inside a safe, checked boundary.
\end{description}

\end{document}
